
\documentclass{article}
\usepackage{colortbl}
\usepackage{makecell}
\usepackage{multirow}
\usepackage{supertabular}

\begin{document}

\newcounter{utterance}

\centering \large Interaction Transcript for game `cladder', experiment `full\_v1.5\_default', episode 5313 with refineonly{-}llama{-}3.1{-}8b{-}instruct.
\vspace{24pt}

{ \footnotesize  \setcounter{utterance}{1}
\setlength{\tabcolsep}{0pt}
\begin{supertabular}{c@{$\;$}|p{.15\linewidth}@{}p{.15\linewidth}p{.15\linewidth}p{.15\linewidth}p{.15\linewidth}p{.15\linewidth}}
    \# & \multicolumn{2}{c}{Player} && \multicolumn{2}{c}{Game Master} \\
    \hline

    \theutterance \stepcounter{utterance}  
    & & & \multicolumn{4}{p{0.6\linewidth}}{
        \cellcolor[rgb]{0.9,0.9,0.9}{
            \makecell[{{p{\linewidth}}}]{
                \texttt{\tiny{[P1$\langle$GM]}}
                \texttt{You are an expert in causal inference. The following question is not a typical commonsense query, but rather a meticulously designed question created by a professor specializing in causal inference, intended to assess the students' mastery of the course content.} \\
\texttt{Imagine a self{-}contained, hypothetical world with only the following conditions, and without any unmentioned factors or causal relationships: Liking spicy food has a direct effect on occupation and salary. Skill has a direct effect on occupation and salary. Occupation has a direct effect on salary. Skill is unobserved. The overall probability of liking spicy food is 2\%. The probability of not liking spicy food and high salary is 45\%. The probability of liking spicy food and high salary is 1\%. Is the chance of high salary smaller when observing liking spicy food?} \\
\texttt{Start your answer with 'Yes' or 'No', followed by additional reasoning or evidence to support your explanation.} \\
            }
        }
    }
     \\ \\

    \theutterance \stepcounter{utterance}  
    & & \multicolumn{4}{p{0.6\linewidth}}{
        \cellcolor[rgb]{0.95,0.95,0.95}{
            \makecell[{{p{\linewidth}}}]{
                \texttt{\tiny{[P1$\leftrightarrow$P]}}
                \texttt{{[}DRAFT{]}} \\
\texttt{No.} \\
\\ 
\texttt{To answer this question, we need to apply the principles of causal inference, specifically the concept of confounding variables and the causal diagram. } \\
\\ 
\texttt{Given the conditions, we have the following causal diagram:} \\
\\ 
\texttt{{-} Liking spicy food (S) → Occupation (O)} \\
\texttt{{-} Skill (X) → Occupation (O)} \\
\texttt{{-} Occupation (O) → Salary (Y)} \\
\texttt{{-} Skill (X) is unobserved.} \\
\\ 
\texttt{We are asked to compare the probability of high salary (Y) when observing liking spicy food (S) versus not observing it. } \\
\\ 
\texttt{To do this, we can use the backdoor criterion, which states that a path from the cause (S) to the effect (Y) is blocked if we condition on all the variables that block the path. In this case, the path from S to Y is blocked by O.} \\
\\ 
\texttt{However, we also have a confounding variable, Skill (X), which affects both S and Y. Since X is unobserved, we cannot condition on it directly.} \\
\\ 
\texttt{To resolve this, we can use the concept of conditional probability and the law of total probability. } \\
\\ 
\texttt{Let's denote the probability of high salary as P(Y=1) and the probability of not liking spicy food as P(S=0). We are given that P(S=0) = 0.98 and P(Y=1|S=0) = 0.45.} \\
\\ 
\texttt{We are also given that P(S=1) = 0.02 and P(Y=1|S=1) = 0.01.} \\
\\ 
\texttt{Using the law of total probability, we can calculate the overall probability of high salary as:} \\
\\ 
\texttt{P(Y=1) = P(Y=1|S=0) \textbackslash{}* P(S=0) + P(Y=1|S=1) \textbackslash{}* P(S=1)} \\
\texttt{= 0.45 \textbackslash{}* 0.98 + 0.01 \textbackslash{}* 0.02} \\
\texttt{= 0.44 + 0.0002} \\
\texttt{= 0.44} \\
\\ 
\texttt{Now, let's calculate the probability of high salary when observing liking spicy food (S=1):} \\
\\ 
\texttt{P(Y=1|S=1) = 0.01} \\
\\ 
\texttt{However, this is not the probability of high salary when we observe S=1. To get the correct probability, we need to condition on the fact that S=1. } \\
\\ 
\texttt{Since S=1 is a rare event (P(S=1) = 0.02), we can assume that the probability of high salary when S=1 is approximately equal to the probability of high salary when S=1 and O is high (i.e., the probability of high salary when S=1 and O is high is approximately equal to the probability of high salary when S=1).} \\
\\ 
\texttt{Let's denote the probability of high salary when S=1 and O is high as P(Y=1|S=1,O=1). We can use the law of total probability} \\
            }
        }
    }
    & \\ \\

    \theutterance \stepcounter{utterance}  
    & \multicolumn{4}{p{0.6\linewidth}}{
        \cellcolor[rgb]{0.561,0.737,0.561}{
            \makecell[{{p{\linewidth}}}]{
                \texttt{\tiny{[P1$\rangle$GM]}}
                \texttt{No.} \\
\\ 
\texttt{To answer this question, we need to apply the principles of causal inference, specifically the concept of confounding variables and the causal diagram.} \\
\\ 
\texttt{Given the conditions, we have the following causal diagram:} \\
\\ 
\texttt{{-} Liking spicy food (S) → Occupation (O)} \\
\texttt{{-} Skill (X) → Occupation (O)} \\
\texttt{{-} Occupation (O) → Salary (Y)} \\
\texttt{{-} Skill (X) is unobserved.} \\
\\ 
\texttt{We are asked to compare the probability of high salary (Y) when observing liking spicy food (S) versus not observing it.} \\
\\ 
\texttt{To do this, we can use the backdoor criterion, which states that a path from the cause (S) to the effect (Y) is blocked if we condition on all the variables that block the path. In this case, the path from S to Y is blocked by O.} \\
\\ 
\texttt{However, we also have a confounding variable, Skill (X), which affects both S and Y. Since X is unobserved, we cannot condition on it directly.} \\
\\ 
\texttt{To resolve this, we can use the concept of conditional probability and the law of total probability.} \\
\\ 
\texttt{Let's denote the probability of high salary as P(Y=1) and the probability of not liking spicy food as P(S=0). We are given that P(S=0) = 0.98 and P(Y=1|S=0) = 0.45.} \\
\\ 
\texttt{We are also given that P(S=1) = 0.02 and P(Y=1|S=1) = 0.01.} \\
\\ 
\texttt{Using the law of total probability, we can calculate the overall probability of high salary as:} \\
\\ 
\texttt{P(Y=1) = P(Y=1|S=0) \textbackslash{}* P(S=0) + P(Y=1|S=1) \textbackslash{}* P(S=1)} \\
\texttt{= 0.45 \textbackslash{}* 0.98 + 0.01 \textbackslash{}* 0.02} \\
\texttt{= 0.44 + 0.0002} \\
\texttt{= 0.44} \\
\\ 
\texttt{Now, let's calculate the probability of high salary when observing liking spicy food (S=1):} \\
\\ 
\texttt{P(Y=1|S=1) = 0.01} \\
\\ 
\texttt{However, this is not the probability of high salary when we observe S=1. To get the correct probability, we need to condition on the fact that S=1.} \\
\\ 
\texttt{Since S=1 is a rare event (P(S=1) = 0.02), we can assume that the probability of high salary when S=1 is approximately equal to the probability of high salary when S=1 and O is high (i.e., the probability of high salary when S=1 and O is high is approximately equal to the probability of high salary when S=1).} \\
\\ 
\texttt{Let's denote the probability of high salary when S=1 and O is high as P(Y=1|S=1,O=1). We can use the law of total probability to calculate this probability} \\
            }
        }
    }
    & & \\ \\

    \theutterance \stepcounter{utterance}  
    & & & \multicolumn{2}{p{0.3\linewidth}}{
        \cellcolor[rgb]{0.9,0.9,0.9}{
            \makecell[{{p{\linewidth}}}]{
                \texttt{\tiny{[GM$|$GM]}}
                \texttt{No} \\
            }
        }
    }
    & & \\ \\

    \theutterance \stepcounter{utterance}  
    & & & \multicolumn{2}{p{0.3\linewidth}}{
        \cellcolor[rgb]{0.9,0.9,0.9}{
            \makecell[{{p{\linewidth}}}]{
                \texttt{\tiny{[GM$|$GM]}}
                \texttt{yes} \\
            }
        }
    }
    & & \\ \\

    \theutterance \stepcounter{utterance}  
    & & & \multicolumn{2}{p{0.3\linewidth}}{
        \cellcolor[rgb]{0.9,0.9,0.9}{
            \makecell[{{p{\linewidth}}}]{
                \texttt{\tiny{[GM$|$GM]}}
                \texttt{game\_result = LOSE} \\
            }
        }
    }
    & & \\ \\

\end{supertabular}
}

\end{document}
